\chapter{Giới Thiệu}

\section{Bối cảnh}

Đồ thị là một mô hình trung tâm của thuật toán ứng dụng. Trong giao thông, nút
biểu diễn giao lộ hoặc điểm trên đường, cạnh biểu diễn đoạn đường và trọng số
biểu diễn chi phí di chuyển. Cùng một cấu trúc đó còn mô tả được những bài toán
động như tìm kiếm, truy đuổi và ra quyết định khi thông tin không đầy đủ. Đồ án
khai thác điểm chung này để xây dựng một ứng dụng vừa có giá trị thực hành, vừa
cho phép quan sát trực tiếp hành vi của nhiều lớp thuật toán.

Phần lõi của đề tài là tối ưu lộ trình đi qua nhiều địa điểm tại Hà Nội. Nếu chỉ
cần đi từ một nguồn tới một đích, bài toán có thể giải bằng đường đi ngắn nhất;
đây cũng là nền của các hệ thống định tuyến quy mô lớn \cite{delling2009}.
Khi cần chọn cả thứ tự ghé nhiều điểm, không gian phương án tăng theo giai thừa;
vì vậy hệ thống kết hợp Dijkstra với quy hoạch động Held--Karp. Phần mở rộng
``Bắt tôi nếu có thể'' biến mạng đường thành một sân pursuit--evasion theo lượt.
Ở đây khó khăn không chỉ nằm ở tìm đường, mà còn ở việc ra quyết định dựa trên
một tập vị trí khả dĩ thay vì biết chính xác đối thủ đang ở đâu.

\section{Bài toán thực tế}

Người dùng cần chọn một điểm xuất phát và từ một đến mười điểm cần ghé. Hệ thống
phải trả về thứ tự ghé, tuyến đường chi tiết, tổng quãng đường và thời gian ước
tính. Ba dạng kết thúc được hỗ trợ:
\begin{itemize}
  \item \textbf{Tuyến mở}: điểm kết thúc do thuật toán tự chọn;
  \item \textbf{Điểm cuối cố định}: người dùng chỉ định điểm cuối;
  \item \textbf{Chu trình}: quay lại điểm xuất phát.
\end{itemize}

Trên cùng giao diện, người dùng có thể chuyển sang trò chơi truy tìm. Người trốn
di chuyển trước với ngân sách tính theo mét. Hai đội máy di chuyển sau, chỉ dựa
trên quan sát công khai. Người trốn thắng khi tới một lối thoát chưa bị chiếm;
máy thắng nếu bắt được người trốn hoặc người trốn hết số lượt. Vị trí thật của
người trốn chỉ được tiết lộ theo chu kỳ cấu hình.

\section{Câu hỏi và mục tiêu}

Đồ án trả lời bốn câu hỏi kỹ thuật:
\begin{enumerate}
  \item Làm thế nào chuyển mạng đường có hướng thành ma trận chi phí giữa các
  điểm quan tâm mà vẫn khôi phục được tuyến đường thật?
  \item Làm thế nào tìm thứ tự ghé tối ưu chính xác trong giới hạn nhỏ, đồng thời
  có baseline heuristic để đối chiếu?
  \item Làm thế nào mô hình hóa trạng thái ẩn của người trốn để chính sách máy
  không vô tình đọc thông tin bí mật?
  \item Làm thế nào kiểm thử các bất biến thuật toán và quy tắc kết thúc thay vì
  chỉ kiểm thử giao diện?
\end{enumerate}

Mục tiêu sản phẩm là một ứng dụng chạy cục bộ, không cần gọi dịch vụ định tuyến
trong luồng tối ưu thông thường; cung cấp phản hồi trực quan trên bản đồ; và duy
trì được trạng thái khi chuyển qua lại giữa hai chế độ. Mục tiêu học thuật là
trình bày rõ mô hình toán, giả định, độ phức tạp, điều kiện đúng và giới hạn của
từng thuật toán.

\section{Phạm vi và giả định}

\begin{table}[H]
\centering
\caption{Phạm vi được thực hiện và nội dung không tuyên bố}
\begin{tabularx}{\linewidth}{Y Y}
\toprule
\textbf{Trong phạm vi} & \textbf{Ngoài phạm vi hoặc chưa được chứng minh} \\
\midrule
Graph ô tô Hà Nội đóng gói cục bộ; tối đa 10 điểm ghé & Dữ liệu giao thông thời
gian thực, ùn tắc hoặc đóng đường phát sinh tức thời \\
Tối ưu chính xác thứ tự ghé theo ma trận chi phí tĩnh & Tối ưu toàn cục cho số
điểm lớn tùy ý hoặc bài toán VRP nhiều xe \cite{toth2014} \\
Trò chơi theo lượt, bắt tại nút cuối bước, tối đa hai đội máy & Va chạm giữa cạnh,
chuyển động liên tục hoặc cân bằng Nash \\
So sánh ba heuristic trên bot và seed cố định & Chứng minh chiến thuật máy tối ưu
hoặc tổng quát hóa cho người chơi thật \\
\bottomrule
\end{tabularx}
\end{table}

Chi phí thời gian trong graph là ước lượng tĩnh được lưu trên cạnh. Vì vậy việc
chọn mục tiêu thời gian không đồng nghĩa với định tuyến phụ thuộc thời gian.
Tương tự, graph trò chơi OSM đi kèm được trích từ một nguồn tải dở và điều chỉnh
hai chiều để phục vụ mô phỏng; nó không được dùng làm hướng dẫn giao thông.

\section{Đóng góp của đồ án}

Các kết quả chính của đồ án gồm:
\begin{itemize}
  \item Một pipeline hoàn chỉnh từ điểm địa lý tới đường đi chi tiết: ánh xạ nút,
  Dijkstra nhiều đích, ma trận chi phí, Held--Karp và khôi phục tuyến.
  \item Một giao diện hợp nhất hai chế độ, giữ độc lập trạng thái lộ trình và ván
  chơi khi người dùng chuyển màn hình.
  \item Một mô hình pursuit--evasion thông tin ẩn với kiểu \texttt{Observation}
  không chứa vị trí hoặc lịch sử bí mật của người trốn.
  \item Ba chính sách truy tìm có giải thích, benchmark 360 ván và bộ 112 kiểm
  thử bao phủ thuật toán, dữ liệu, quy tắc biên và tích hợp giao diện.
\end{itemize}

